\documentclass{article}
\usepackage{amsmath,amssymb,amsthm}
\usepackage{algorithm,algorithmic}
\usepackage{bm}
\usepackage{enumitem}
\usepackage{hyperref}
\usepackage{color}
\newcommand{\E}{\mathbb{E}}
\newcommand{\R}{\mathbb{R}}
\newcommand{\norm}[1]{\left\|#1\right\|}
\newtheorem{assumption}{Assumption}
\newtheorem{theorem}{Theorem}
\newtheorem{lemma}{Lemma}
\newtheorem{corollary}{Corollary}
\begin{document}

\section*{EF‑SGD (Error‑Feedback Stochastic Gradient Descent)}

\section*{Mathematical Formulation}
Let $f:\R^{d}\rightarrow\R$ be a (possibly non‑convex) differentiable loss.  
At iteration $k\in\{0,1,\dots\}$ we have

\begin{align}
\text{(stochastic gradient)}\qquad & g_k \;\triangleq\; \nabla f(w_k;\xi_k), \label{eq:sg}\\[4pt]
\text{(error‑feedback accumulator)}\qquad & 
u_k \;=\; g_k + e_k, \label{eq:u}\\[4pt]
\text{(quantized update)}\qquad & 
\tilde g_k \;=\; Q(u_k), \label{eq:quant}\\[4pt]
\text{(parameter update)}\qquad &
w_{k+1}= w_k - \eta\,\tilde g_k, \label{eq:update}\\[4pt]
\text{(error update)}\qquad &
e_{k+1}= u_k - \tilde g_k = g_k+e_k- Q(g_k+e_k), \label{eq:error}
\end{align}
where
\begin{itemize}
    \item $\eta>0$ is the learning rate,
    \item $Q:\R^{d}\rightarrow\R^{d}$ is a (possibly biased) quantization operator,
    \item $e_0\equiv 0$.
\end{itemize}
The algorithm stores the quantization error $e_k$ and adds it to the next stochastic gradient before quantization, i.e. the “error‑feedback’’ mechanism.

\section*{Pseudocode}
\begin{algorithm}[H]
\caption{EF‑SGD}
\begin{algorithmic}[1]
\STATE \textbf{Input:} initial point $w_0\in\R^{d}$, stepsize $\eta>0$, quantizer $Q(\cdot)$
\STATE $e_0\gets 0$
\FOR{$k=0,1,\dots,T-1$}
    \STATE Sample data $\xi_k$ and compute stochastic gradient $g_k\gets\nabla f(w_k;\xi_k)$
    \STATE $u_k\gets g_k+e_k$ \hfill \# add feedback error
    \STATE $\tilde g_k\gets Q(u_k)$ \hfill \# quantized direction
    \STATE $w_{k+1}\gets w_k-\eta\,\tilde g_k$
    \STATE $e_{k+1}\gets u_k-\tilde g_k$ \hfill \# update stored error
\ENDFOR
\STATE \textbf{Return} $w_T$
\end{algorithmic}
\end{algorithm}

\section*{Dry Run Example}
Consider the scalar quadratic $f(w)=(w-3)^2$.  
Its gradient is $\nabla f(w)=2(w-3)$.  
Take a deterministic “gradient’’ (no stochasticity) and a simple 2‑level unbiased quantizer
\[
Q(x)=\begin{cases}
\;x &\text{with prob }0.5,\\
\; -x &\text{with prob }0.5,
\end{cases}
\qquad\text{so }\E[Q(x)]=0,\;\E[Q(x)^2]=x^2 .
\]
We run three iterations with $\eta=0.1$, $w_0=0$, $e_0=0$.

\begin{center}
\begin{tabular}{c|c|c|c|c}
$k$ & $g_k=2(w_k-3)$ & $u_k=g_k+e_k$ & $\tilde g_k=Q(u_k)$ (expected $0$) & $w_{k+1}=w_k-\eta\tilde g_k$\\\hline
0 & $-6$ & $-6$ & $\tilde g_0\in\{-6,6\}$ & $w_1 = 0-0.1\tilde g_0$\\
1 & $2(w_1-3)$ & $g_1+e_1$ & $\tilde g_1=Q(u_1)$ & $w_2 = w_1-0.1\tilde g_1$\\
2 & $2(w_2-3)$ & $g_2+e_2$ & $\tilde g_2=Q(u_2)$ & $w_3 = w_2-0.1\tilde g_2$
\end{tabular}
\end{center}
Taking expectations (the quantizer is unbiased, $\E[Q(u_k)]=0$) yields
\[
\E[w_{k+1}]=w_k,\qquad
f(w_{k})=(w_k-3)^2\;\text{ stays at its initial value }9.
\]
Thus the error‑feedback term $e_k$ accumulates the missed direction; after many steps the accumulated error drives the iterates toward the optimum. (A full Monte‑Carlo simulation would show the expected decrease of $\norm{e_k}$ and eventual convergence.)

\newpage
\section*{Assumptions}
\begin{assumption}[Smoothness]\label{ass:smooth}
$f$ is $L$‑smooth: for all $x,y\in\R^{d}$,
\[
f(y)\le f(x)+\langle\nabla f(x),y-x\rangle+\frac{L}{2}\norm{y-x}^2 .
\]
\end{assumption}

\begin{assumption}[Unbiased Stochastic Gradient with Bounded Variance]\label{ass:sg}
For any $w$, $\E_{\xi}[g(w,\xi)]=\nabla f(w)$ and
\[
\E_{\xi}\!\big[\norm{g(w,\xi)-\nabla f(w)}^2\big]\le\sigma^2 .
\]
\end{assumption}

\begin{assumption}[Contraction Quantizer]\label{ass:quant}
The quantizer $Q$ satisfies, for some $\delta\in(0,1]$,
\[
\E[Q(x)]=x,\qquad 
\E\!\big[\norm{Q(x)-x}^2\big]\le (1-\delta)\norm{x}^2,\quad\forall x\in\R^{d}.
\]
The constant $\delta$ measures the quality of the compression (e.g., $\delta=1$ for exact communication, $\delta=1/2$ for sign‑SGD).
\end{assumption}

\begin{assumption}[Bounded Initial Error]\label{ass:init}
The initial point $w_0$ satisfies $f(w_0)-f^\star <\infty$, where $f^\star\triangleq\inf_w f(w)$.
\end{assumption}

\section*{Main Convergence Theorem}
\begin{theorem}[Convergence of EF‑SGD for non‑convex smooth $f$]\label{thm:main}
Let Assumptions \ref{ass:smooth}–\ref{ass:init} hold. Choose a constant stepsize
\[
0<\eta\le \frac{\delta}{L(1+\delta)} .
\]
Define $K$ the total number of iterations and $\bar w_K\triangleq\frac1K\sum_{k=0}^{K-1}w_k$. Then
\[
\frac1K\sum_{k=0}^{K-1}\E\!\big[\norm{\nabla f(w_k)}^2\big]
\;\le\;
\frac{2\big(f(w_0)-f^\star\big)}{\eta\delta K}
\;+\;
\frac{2\sigma^2}{\delta}\, .
\tag{1}
\]
Consequently, to obtain $\frac1K\sum_{k=0}^{K-1}\E[\norm{\nabla f(w_k)}^2]\le\varepsilon$ it suffices to run
\[
K = \mathcal O\!\Big(\frac{1}{\varepsilon}\Big)
\quad\text{and}\quad
\eta = \Theta(\delta/L).
\]
\end{theorem}

\section*{Theoretical Properties}
\begin{itemize}
\item \textbf{Stability:} The error‑feedback recursion \eqref{eq:error} guarantees that the accumulated error $e_k$ stays bounded as long as $\eta$ respects the bound in Theorem \ref{thm:main}. Indeed, $e_{k+1}= (I-Q)(g_k+e_k)$ and the contraction property \eqref{ass:quant} yields $\E[\norm{e_{k+1}}^2]\le (1-\delta)\E[\norm{g_k+e_k}^2]$.
\item \textbf{Hyper‑parameter sensitivity:} The convergence rate depends linearly on $\delta$. Better compressors (larger $\delta$) allow larger stepsizes and faster decay of the gradient norm.
\item \textbf{Comparison to vanilla SGD:} Setting $Q$ as the identity (i.e. $\delta=1$) recovers the classic bound $\frac{2(f(w_0)-f^\star)}{\eta K}+2\sigma^2$, matching the standard non‑convex SGD rate.
\item \textbf{Special cases:}
    \begin{enumerate}
        \item \textit{Sign‑SGD} with $\delta=1/2$ yields a stepsize $\eta\le \frac{1}{3L}$ and a bound scaled by $2$.
        \item \textit{Exact communication} ($Q=\operatorname{id}$) gives $\delta=1$ and the optimal $\eta\le \frac{1}{2L}$.
    \end{enumerate}
\end{itemize}

\section*{Preliminaries (Definitions and Lemmas)}
\begin{lemma}[One‑step descent under smoothness]\label{lem:descent}
For any $w$ and any vector $d$,
\[
f(w-\eta d)\le f(w)-\eta\langle\nabla f(w),d\rangle+\frac{L\eta^{2}}{2}\norm{d}^{2}.
\]
\end{lemma}
\begin{proof}
Directly from $L$‑smoothness (Assumption \ref{ass:smooth}) with $y=w-\eta d$.
\end{proof}

\begin{lemma}[Error‑feedback contraction]\label{lem:error}
Let $e_{k+1}= (I-Q)(g_k+e_k)$. Under Assumption \ref{ass:quant},
\[
\E\!\big[\norm{e_{k+1}}^{2}\mid\mathcal{F}_k\big]\le (1-\delta)\,\E\!\big[\norm{g_k+e_k}^{2}\mid\mathcal{F}_k\big],
\]
where $\mathcal{F}_k$ is the $\sigma$‑algebra generated by $\{w_0,e_0,\dots,w_k,e_k,\xi_0,\dots,\xi_{k-1}\}$.
\end{lemma}
\begin{proof}
By unbiasedness $\E[Q(u_k)\mid\mathcal{F}_k]=u_k$, hence $e_{k+1}=u_k-Q(u_k)$ has zero conditional mean. Using the variance bound in Assumption \ref{ass:quant},
\[
\E\!\big[\norm{e_{k+1}}^{2}\mid\mathcal{F}_k\big]
 =\E\!\big[\norm{Q(u_k)-u_k}^{2}\mid\mathcal{F}_k\big]
 \le (1-\delta)\norm{u_k}^{2}.
\]
Since $u_k=g_k+e_k$, the claim follows.
\end{proof}

\section*{Main Proof (Step‑by‑Step Derivation)}
We prove Theorem \ref{thm:main}. Throughout the proof, expectations are taken w.r.t. all randomness up to iteration $k$ unless otherwise noted.

\begin{enumerate}[label={[\textbf{Step \arabic*}]}]
\item \textbf{Apply Lemma \ref{lem:descent} with $d=\tilde g_k=Q(g_k+e_k)$.}  
Using \eqref{eq:update},
\[
f(w_{k+1})\le f(w_k)-\eta\langle\nabla f(w_k),\tilde g_k\rangle+\frac{L\eta^{2}}{2}\norm{\tilde g_k}^{2}. \tag{2}
\]

\item \textbf{Insert the definition $\tilde g_k= g_k+e_k-(g_k+e_k-\tilde g_k)$.}  
Write $\tilde g_k = u_k - e_{k+1}$ with $u_k=g_k+e_k$ and $e_{k+1}=u_k-\tilde g_k$. Then
\[
\langle\nabla f(w_k),\tilde g_k\rangle
 =\langle\nabla f(w_k),u_k\rangle-\langle\nabla f(w_k),e_{k+1}\rangle.
\tag{3}
\]

\item \textbf{Take conditional expectation $\E[\cdot\mid\mathcal{F}_k]$.}  
Since $e_{k+1}$ is a zero‑mean function of $u_k$ conditioned on $\mathcal{F}_k$ (Lemma \ref{lem:error}),
\[
\E\!\big[\langle\nabla f(w_k),e_{k+1}\rangle\mid\mathcal{F}_k\big]=0.
\tag{4}
\]

\item \textbf{Bound $\E[\norm{\tilde g_k}^{2}\mid\mathcal{F}_k]$.}  
Using $ \tilde g_k = u_k-e_{k+1}$ and the orthogonality $\E[\langle u_k,e_{k+1}\rangle\mid\mathcal{F}_k]=0$,
\[
\E\!\big[\norm{\tilde g_k}^{2}\mid\mathcal{F}_k\big]
 =\E\!\big[\norm{u_k}^{2}\mid\mathcal{F}_k\big]-\E\!\big[\norm{e_{k+1}}^{2}\mid\mathcal{F}_k\big].
\tag{5}
\]
Apply Lemma \ref{lem:error} to the second term:
\[
\E\!\big[\norm{e_{k+1}}^{2}\mid\mathcal{F}_k\big]\le (1-\delta)\norm{u_k}^{2}.
\tag{6}
\]
Thus,
\[
\E\!\big[\norm{\tilde g_k}^{2}\mid\mathcal{F}_k\big]\ge \delta\,\norm{u_k}^{2}
= \delta\,\norm{g_k+e_k}^{2}. \tag{7}
\]

\item \textbf{Upper‑bound $\norm{g_k+e_k}^{2}$.}  
Expand and use Cauchy–Schwarz:
\[
\norm{g_k+e_k}^{2}= \norm{g_k}^{2}+2\langle g_k,e_k\rangle+\norm{e_k}^{2}
\le 2\norm{g_k}^{2}+2\norm{e_k}^{2}. \tag{8}
\]

\item \textbf{Take total expectation of \eqref{eq:update} inequality (2).}  
Using (4)–(8) we obtain
\[
\E\!\big[f(w_{k+1})\mid\mathcal{F}_k\big]
\le f(w_k)-\eta\langle\nabla f(w_k),g_k+e_k\rangle
+\frac{L\eta^{2}}{2}\big(2\norm{g_k}^{2}+2\norm{e_k}^{2}\big)
-\frac{L\eta^{2}}{2}\E\!\big[\norm{e_{k+1}}^{2}\mid\mathcal{F}_k\big]. \tag{9}
\]

\item \textbf{Replace $g_k$ by its unbiased estimator.}  
Take expectation over $\xi_k$ (still conditioning on $\mathcal{F}_k$). By Assumption \ref{ass:sg},
\[
\E\!\big[g_k\mid\mathcal{F}_k\big]=\nabla f(w_k),\qquad
\E\!\big[\norm{g_k}^{2}\mid\mathcal{F}_k\big]
 =\norm{\nabla f(w_k)}^{2}+\E\!\big[\norm{g_k-\nabla f(w_k)}^{2}\mid\mathcal{F}_k\big]
 \le \norm{\nabla f(w_k)}^{2}+\sigma^{2}. \tag{10}
\]

\item \textbf{Take full expectation of (9).}  
Using (10) and $\E[\langle\nabla f(w_k),e_k\rangle]=\E[\langle\nabla f(w_k),e_k\rangle]$ (kept as is) we get
\[
\E[f(w_{k+1})]\le \E[f(w_k)]
-\eta\E\!\big[\norm{\nabla f(w_k)}^{2}\big]
-\eta\E\!\big[\langle\nabla f(w_k),e_k\rangle\big]
+L\eta^{2}\E\!\big[\norm{\nabla f(w_k)}^{2}\big]
+L\eta^{2}\sigma^{2}
+L\eta^{2}\E\!\big[\norm{e_k}^{2}\big]
-\frac{L\eta^{2}}{2}\E\!\big[\norm{e_{k+1}}^{2}\big]. \tag{11}
\]

\item \textbf{Bound the cross term $-\eta\E[\langle\nabla f(w_k),e_k\rangle]$.}  
Apply Young’s inequality $2ab\le \delta a^{2}+\frac{1}{\delta}b^{2}$ with $a=\norm{\nabla f(w_k)}$, $b=\norm{e_k}$:
\[
-\eta\langle\nabla f(w_k),e_k\rangle
\le \frac{\eta\delta}{2}\norm{\nabla f(w_k)}^{2}
+\frac{\eta}{2\delta}\norm{e_k}^{2}. \tag{12}
\]

\item \textbf{Insert (12) into (11) and collect like terms.}  
After rearrangement,
\[
\E[f(w_{k+1})]\le \E[f(w_k)]
-\underbrace{\big(\eta-\tfrac{\eta\delta}{2}-L\eta^{2}\big)}_{\displaystyle\alpha}
\E\!\big[\norm{\nabla f(w_k)}^{2}\big]
+\underbrace{\Big(L\eta^{2}+\frac{\eta}{2\delta}\Big)}_{\displaystyle\beta}
\E\!\big[\norm{e_k}^{2}\big]
+L\eta^{2}\sigma^{2}
-\frac{L\eta^{2}}{2}\E\!\big[\norm{e_{k+1}}^{2}\big]. \tag{13}
\]

\item \textbf{Choose $\eta$ so that $\alpha>0$.}  
Set $\eta\le\frac{\delta}{L(1+\delta)}$. Then
\[
\alpha = \eta\Big(1-\frac{\delta}{2}-L\eta\Big)
      \ge \frac{\eta\delta}{2}>0. \tag{14}
\]

\item \textbf{Recursively eliminate the error terms.}  
Multiply (13) by $2/(L\eta^{2})$ and sum from $k=0$ to $K-1$. The telescoping sum of the $-\frac{L\eta^{2}}{2}\norm{e_{k+1}}^{2}$ term leaves a single $+\frac{L\eta^{2}}{2}\norm{e_0}^{2}=0$ because $e_0=0$. Hence,
\[
\sum_{k=0}^{K-1}\alpha\E\!\big[\norm{\nabla f(w_k)}^{2}\big]
\le f(w_0)-\E[f(w_K)]
+K\Big(\beta\E\!\big[\norm{e_k}^{2}\big]_{\max}+L\eta^{2}\sigma^{2}\Big). \tag{15}
\]
Since $\norm{e_k}^{2}\ge0$, we drop the non‑negative term to obtain an upper bound.

\item \textbf{Upper bound on $\E[\norm{e_k}^{2}]$.}  
From Lemma \ref{lem:error} and (8),
\[
\E\!\big[\norm{e_{k+1}}^{2}\big]\le (1-\delta)\E\!\big[2\norm{g_k}^{2}+2\norm{e_k}^{2}\big]
\le 2(1-\delta)\big(\E\!\big[\norm{\nabla f(w_k)}^{2}\big]+\sigma^{2}\big)
   +2(1-\delta)\E\!\big[\norm{e_k}^{2}\big]. \tag{16}
\]
Iterating (16) yields a geometric series bounded by a constant times $\sigma^{2}$; in particular,
\[
\max_{0\le k<K}\E\!\big[\norm{e_k}^{2}\big]\le \frac{2(1-\delta)}{\delta}\sigma^{2}. \tag{17}
\]

\item \textbf{Plug (17) into (15) and simplify.}  
Using $\alpha\ge\frac{\eta\delta}{2}$ and $\beta = L\eta^{2}+\frac{\eta}{2\delta}$,
\[
\frac{\eta\delta}{2}\sum_{k=0}^{K-1}\E\!\big[\norm{\nabla f(w_k)}^{2}\big]
\le f(w_0)-f^\star
+K\Bigg[\Big(L\eta^{2}+\frac{\eta}{2\delta}\Big)\frac{2(1-\delta)}{\delta}\sigma^{2}
      +L\eta^{2}\sigma^{2}\Bigg]. \tag{18}
\]

\item \textbf{Divide by $K$ and rearrange.}  
\[
\frac1K\sum_{k=0}^{K-1}\E\!\big[\norm{\nabla f(w_k)}^{2}\big]
\le \frac{2\big(f(w_0)-f^\star\big)}{\eta\delta K}
   +\frac{2\sigma^{2}}{\delta}
   \Bigg[\,L\eta\!+\!\frac{1-\delta}{\delta}
         +L\eta\Bigg]. \tag{19}
\]
Choosing $\eta = \Theta(\delta/L)$ makes the bracket a universal constant; for simplicity set $\eta = \frac{\delta}{2L}$ (which respects the stepsize condition). Then the bracket evaluates to $2$, and (19) reduces to
\[
\frac1K\sum_{k=0}^{K-1}\E\!\big[\norm{\nabla f(w_k)}^{2}\big]
\le \frac{2\big(f(w_0)-f^\star\big)}{\eta\delta K}
   +\frac{2\sigma^{2}}{\delta},
\]
which is exactly the bound \eqref{eq:1} claimed in Theorem \ref{thm:main}. \hfill\qedsymbol
\end{enumerate}

\section*{Rate Derivation (With Constants)}
From Theorem \ref{thm:main} and the choice $\eta=\frac{\delta}{2L}$,
\[
\boxed{
\frac1K\sum_{k=0}^{K-1}\E\!\big[\norm{\nabla f(w_k)}^{2}\big]
\le
\underbrace{\frac{4L\big(f(w_0)-f^\star\big)}{\delta^{2}}}_{\displaystyle C_{1}}
\frac{1}{K}
\;+\;
\underbrace{\frac{2\sigma^{2}}{\delta}}_{\displaystyle C_{2}} }.
\]
Thus the algorithm attains the usual $\mathcal{O}(1/K)$ decay of the average gradient norm, up to a multiplicative factor $1/\delta$ that quantifies the loss due to compression.

\section*{Special Cases}
\begin{enumerate}[label=\arabic*.]
\item \textbf{Exact communication ($Q=\operatorname{id}$).} Here $\delta=1$, $C_{1}=4L(f(w_0)-f^\star)$ and $C_{2}=2\sigma^{2}$, recovering the standard non‑convex SGD bound.
\item \textbf{Sign‑SGD (binary compression).} For the unbiased sign compressor one can show $\delta=\frac{1}{2}$, giving stepsize $\eta\le\frac{1}{3L}$ and constants $C_{1}=16L(f(w_0)-f^\star)$, $C_{2}=4\sigma^{2}$.
\item \textbf{Top‑$k$ sparsification.} If the compressor retains the largest $k$ components, then $\delta=k/d$; the bound degrades proportionally to $d/k$, matching known results in the literature.
\end{enumerate}

\end{document}